\paragraph{The cost of grid-following wind, and support without a reserve.} Fig.~\ref{fig:B1_frequency} and Table~\ref{tab:B} show a 50 MW (5 percent) load step at $t=1$ s with 40 percent of the load supplied by wind at 9 m/s. Replacing synchronous generation by grid-following MPPT wind (C1) cuts the system inertia constant from 5.88 s (C0) to 3.76 s: the nadir falls from 49.771 to 49.659 Hz and the 500 ms RoCoF rises from 0.198 to 0.305 Hz/s. The MPPT turbine does not react to the event at all, its rotor stays at 0.750 pu. Adding synthetic inertia and droop to the MPPT turbine (C2) lowers the 500 ms RoCoF to 0.122 Hz/s, but the support is borrowed from the rotor alone. The rotor slows from 0.750 to 0.701 pu, leaves the $C_p$ optimum, and aerodynamic power falls (Fig.~\ref{fig:B2_wind_response}). The frequency therefore keeps sagging to a late minimum of 49.815 Hz at 11.2 s after the step and settles at $-0.160$ Hz, worse than no support at all (C1, $-0.148$ Hz). This is the secondary-dip problem of non-deloaded turbines \cite{ullah}, appearing here as a prolonged sag.

\begin{table}[H]\centering\small\setlength{\tabcolsep}{4pt}
\caption{Frequency metrics after the 50 MW load step (40 percent wind, 9 m/s). $t_{nadir}$ is measured from the step; the last column is the extra wind energy delivered in the 10 s after the step.}\label{tab:B}
\begin{tabular}{@{}lccccccc@{}}\toprule
Case & \shortstack{Nadir\\(Hz)} & \shortstack{$t_{nadir}$\\(s)} & \shortstack{RoCoF inst.\\(Hz/s)} & \shortstack{RoCoF 500 ms\\(Hz/s)} & \shortstack{$\Delta f$ at 40 s\\(Hz)} & \shortstack{min $\omega_r$\\(pu)} & \shortstack{Extra energy\\(MJ)}\\\midrule
C0 all-synchronous & 49.771 & 2.14 & 0.213 & 0.198 & $-0.097$ & n/a & 0\\
C1 GFL, MPPT & 49.659 & 2.09 & 0.332 & 0.305 & $-0.148$ & 0.750 & 0\\
C2 GFL+SI+droop, no reserve & 49.815 & 11.20 & 0.332 & 0.122 & $-0.160$ & 0.701 & 167\\
C3 GFL+SI+droop, 10\% reserve & 49.906 & 2.07 & 0.332 & 0.115 & $-0.070$ & 0.831 & 317\\
C4 GFM-VSM, 10\% reserve & 49.906 & 2.03 & 0.162 & 0.116 & $-0.070$ & 0.831 & 317\\\bottomrule
\end{tabular}\end{table}

\paragraph{What the reserve does, and what grid-forming adds.} With the 10 percent overspeed reserve (C3, C4) the rotor slows only from 0.888 to 0.831 pu and stays above the MPPT speed of 0.750 pu. Slowing therefore raises $P_{aero}$ along the overspeed side of Fig.~\ref{fig:A2_power_speed}, as explained with Eq.~(\ref{eq:gfm}): the plant delivers about 37 MW extra (peak 36.9 MW in C3 and 37.2 MW in C4), 317 MJ in 10 s. The nadir is 49.906 Hz, inside the IEGC band of 49.90 to 50.05 Hz, the steady-state deviation is $-0.070$ Hz and there is no secondary dip. GFM (C4) and GFL with the same reserve (C3) reach practically the same nadir (49.906 Hz in both), because the reserve sets the nadir. What GFM adds is at the step. Its instantaneous RoCoF is half that of C3 (0.162 against 0.332 Hz/s; centre-of-inertia value 0.148 against 0.268 Hz/s), because it acts as a voltage source and needs no frequency measurement, and it adds real inertia: $H_{sys}$ rises from 3.76 to 8.51 s when the virtual inertia of the GFM plant is counted. Over a 500 ms window the two are similar (about 0.116 Hz/s), because our GFL frequency-measurement lag of 50 ms is optimistic. Step 2 will quantify realistic measurement windows. The decisive GFM advantage is weak-grid stability (Task C below).

\paragraph{Penetration and weak connection.} Fig.~\ref{fig:B4_penetration} sweeps the wind share from 10 to 60 percent. The GFL-MPPT nadir worsens from 49.751 to 49.547 Hz, while the GFM nadir improves from 49.811 to 49.928 Hz, because every added GFM turbine adds both virtual inertia and reserve. The support-without-reserve case C2 is not monotonic (49.807 Hz at 10 percent, 49.827 Hz at 30 percent, 49.685 Hz at 60 percent) because it has no reserve to draw on. In Fig.~\ref{fig:B5_weak_connection} the connection of the GFM plant is weakened from $X_{23}=0.1$ to $0.8$ pu, so that the SCR at the wind bus falls from 3.01 to 1.00. The nadir stays at 49.906 Hz, the 500 ms RoCoF rises modestly from 0.112 to 0.148 Hz/s, and the damping ratio of the synchronous-generator/VSM electromechanical mode rises from 0.250 to 0.600 (mode frequency 1.83 to 0.89 Hz).
